% !PRINTED 687
\weberpage{687}
The number $T$ is a function of $t$ which in general cannot be
 determined more precisely; it grows to infinity together with $t$,
but, since it takes only integral values, it changes only
stepwise, and therefore discontinuously. What matters here is
\gesperrt{the limiting value which the ratio $T:t$ approaches as
$t$ grows infinitely} (which in the example cited
 above is $1:m$).

This limiting value can be found by considerations of the kind
we have already used in \S.\,155.

We first form a basis form of the given ideal $\mathfrak{a}$:
\weq{1}{
\alpha = \alpha_1 x_1 + \alpha_2 x_2 + \cdots + \alpha_n x_n,
}
that is, a linear function of the $n$ variables $x_i$ with the property
that, by substituting all integral rational numbers for
$x_i$, there arise from $\alpha$ all the integers of the field $\Omega$ divisible by $\mathfrak{a}$
 (\S.\,146). By $\alpha_{i,1}, \alpha_{i,2}, \ldots, \alpha_{i,n}$ we denote the values conjugate
to $\alpha_i$, and we form the linear substitution for $x_i$:
\weq{2}{
\begin{aligned}
y_1 &= \alpha_{1,1} x_1 + \alpha_{2,1} x_2 + \cdots + \alpha_{n,1} x_n\\
y_2 &= \alpha_{1,2} x_1 + \alpha_{2,2} x_2 + \cdots + \alpha_{n,2} x_n\\
&\ \cdot\ \cdot\ \cdot\ \cdot\ \cdot\ \cdot\ \cdot\ \cdot\ \cdot\ \cdot\ \cdot\ \cdot\\
y_n &= \alpha_{1,n} x_1 + \alpha_{2,n} x_2 + \cdots + \alpha_{n,n} x_n.
\end{aligned}
}
If $A$ is the determinant of this system of equations
\[
A = \Sigma \pm \alpha_{1,1} \alpha_{2,2} \ldots \alpha_{n,n},
\]
then $A^2$ is the discriminant of the system of functions $(\alpha_1, \alpha_2, \ldots, \alpha_n)$,
and therefore, by \S.\,147, equal to the product of the square
of the absolute norm of $\alpha$ and the field discriminant $\Delta$.
Since the absolute norm of $\alpha$ is at the same time the norm of the ideal $\mathfrak{a}$
(\S.\,151), we set
\weq{3}{
A^2 = N(\mathfrak{a})^2 \, \Delta.
}
Now, as in \S.\,155, we may interpret $x_1, x_2, \ldots, x_n$ as
coordinates of a point $(x)$ in a space of $n$ dimensions,
 and we let a point $(x')$ correspond to every point $(x)$,
whose coordinates are determined by
\weq{4}{
x_1' = t^{\frac{1}{n}} x_1, \quad x_2' = t^{\frac{1}{n}} x_2, \quad \ldots, \quad x_n' = t^{\frac{1}{n}} x_n
}
where $t$ denotes an arbitrary positive quantity.
To every region $S$ of points $(x)$ there corresponds a region $S'$ of
points $(x')$. All figures in $S$ are similar to the corresponding figures
in $S'$, and the linear dimensions in $S$ are to
those in $S'$ as $1 : t^{\frac{1}{n}}$.
